\documentclass[11pt,a4paper]{article}
\usepackage[utf8]{inputenc}
\usepackage[T1]{fontenc}
\usepackage{tgheros}
\renewcommand{\familydefault}{\sfdefault}
\usepackage[varqu,scaled=0.92]{zi4}
\usepackage{microtype}
\DisableLigatures[-,<,>]{encoding=T1,family=tt*}
\usepackage{needspace}
\usepackage[a4paper,top=2.5cm,bottom=2.5cm,left=2.5cm,right=2.2cm]{geometry}
\usepackage{graphicx}
\graphicspath{{../../template/}{figures/}}
\usepackage{xcolor}
\usepackage{listings}
\usepackage{upquote}
\usepackage{titlesec}
\usepackage{tocloft}
\usepackage{enumitem}
\usepackage{booktabs}
\usepackage{array}
\usepackage{float}
\usepackage[hidelinks]{hyperref}

\setlength{\parindent}{0pt}
\setlength{\parskip}{0.55em}
\setlist{itemsep=0.3em,topsep=0.3em}
\pagestyle{plain}

\titleformat{\section}{\Large\bfseries}{\thesection}{0.9em}{}
\titleformat{\subsection}{\large\bfseries}{\thesubsection}{0.9em}{}
\titlespacing*{\section}{0pt}{1.6em}{0.7em}
\titlespacing*{\subsection}{0pt}{1.2em}{0.5em}

\renewcommand{\contentsname}{Cuprins}
\renewcommand{\cfttoctitlefont}{\LARGE\bfseries}
\renewcommand{\cftsecfont}{\bfseries}
\renewcommand{\cftsecpagefont}{\bfseries}
\renewcommand{\cftsubsecfont}{\small}
\renewcommand{\cftsubsecpagefont}{\small}
\setlength{\cftbeforesecskip}{0.6em}

\definecolor{codebg}{HTML}{F6F6F6}
\definecolor{coderule}{HTML}{8C8C8C}
\definecolor{codenum}{HTML}{8C8C8C}
\definecolor{codecomment}{HTML}{6A6A6A}

\lstdefinestyle{term}{
  basicstyle=\ttfamily\small,
  backgroundcolor=\color{codebg},
  frame=single,
  rulecolor=\color{coderule},
  framesep=4pt,
  numbers=left,
  numberstyle=\tiny\color{codenum},
  numbersep=9pt,
  xleftmargin=1.4em,
  breaklines=true,
  postbreak=\mbox{\textcolor{codenum}{$\hookrightarrow$}\space},
  columns=fullflexible,
  keepspaces=true,
  showstringspaces=false,
  upquote=true,
  aboveskip=0.9em,
  belowskip=0.9em
}
\lstdefinestyle{ccode}{
  style=term,
  language=C,
  keywordstyle=\bfseries,
  commentstyle=\color{codecomment},
  morekeywords={uint,uint64}
}
\lstset{style=term,captionpos=b}

\usepackage{caption}
\captionsetup[lstlisting]{labelsep=endash,justification=centering,singlelinecheck=false,font=small,skip=7pt}
\captionsetup[table]{labelsep=endash,font=small,skip=6pt}
\renewcommand{\lstlistingname}{Figure}
\AtBeginDocument{\counterwithin{lstlisting}{section}\counterwithin{table}{section}}

\newcommand{\plainsection}[1]{\section*{#1}\addcontentsline{toc}{section}{#1}}
\newcommand{\code}[1]{\texttt{#1}}

\newcommand{\utmtitlepage}{%
\begin{titlepage}
  \centering
  \includegraphics[width=4.3cm]{utm-logo.jpeg}\par
  \vspace{1.3cm}
  {\large\bfseries UNIVERSITATEA TEHNICĂ A MOLDOVEI\par}
  \vspace{0.25cm}
  Facultatea Calculatoare, Informatică și Microelectronică\par
  Departamentul Ingineria Software și Automatică\par
  \vspace{2.6cm}
  {\Large\bfseries LUCRARE DE LABORATOR NR.~\labnumber\par}
  \vspace{0.2cm}
  {\LARGE\bfseries \labtitle\par}
  \vspace{0.45cm}
  Disciplina: Criptografie și Securitate\par
  \ifx\variant\empty\else\vspace{0.2cm}Varianta \variant\par\fi
  \vspace{2.8cm}
  \begin{minipage}[t]{0.46\textwidth}
    \raggedright
    \textbf{A efectuat:}\\
    student gr. \studentgroup\\
    \studentname
  \end{minipage}\hfill
  \begin{minipage}[t]{0.46\textwidth}
    \raggedleft
    \textbf{A verificat:}\\
    conf. univ., dr. Aureliu \textsc{Zgureanu}
  \end{minipage}
  \vfill
  Chișinău — 2026\par
\end{titlepage}}

\usepackage{textcomp}
\usepackage{framed}
\newenvironment{cipherbox}{%
  \def\FrameCommand{\fboxsep=6pt\fboxrule=0.4pt\fcolorbox{coderule}{codebg}}%
  \MakeFramed{\advance\hsize-\width\FrameRestore}%
  \ttfamily\footnotesize\raggedright\frenchspacing\setlength{\parskip}{0pt}}%
  {\par\endMakeFramed}

\newcommand{\tool}{Wolfram Cloud, the online version of Wolfram Mathematica}
\newcommand{\toolshort}{Wolfram Cloud}

\newcommand{\labnumber}{5}
\newcommand{\labtitle}{Public-Key Cryptography: RSA, ElGamal and Diffie-Hellman}
\newcommand{\studentname}{Daniel Nicorici}
\newcommand{\studentgroup}{FAF-243}
\newcommand{\variant}{}
\lstdefinestyle{wl}{style=term,language=Mathematica,keywordstyle=\bfseries}
\lstdefinestyle{py}{style=term,language=Python,keywordstyle=\bfseries}
\begin{document}
\utmtitlepage
\tableofcontents
\newpage

\plainsection{Purpose of the work}

In a public-key cryptosystem every user has two keys: a public key that anyone may use to encrypt
or verify, and a private key that only the owner knows. The purpose of this lab is to generate keys
and to encrypt and decrypt a message with the two classic public-key algorithms, RSA and ElGamal,
and to perform a Diffie-Hellman key exchange that gives Alice and Bob a shared 256-bit key for AES.

The message is my name written as in the handout, surname first: $m = $ \textit{Nicorici Daniel}.
RSA uses a modulus $n$ of 2048 bits. ElGamal and Diffie-Hellman use the 2048-bit prime $p$ and the
generator $g = 2$ given in the handout.

The handout asks for Wolfram|Alpha or Wolfram Mathematica. Typing several 617-digit numbers into
Wolfram|Alpha is not practical, so all three tasks are written as one Wolfram Language program
(the language of Mathematica), using exactly the functions recommended in the handout:
\code{RandomPrime}, \code{RandomInteger}, \code{PowerMod}, \code{Mod}, \code{IntegerDigits}. The
results in this report were produced by running it in \tool. Every step was then checked
independently by a short Python script (section~\ref{sec:verify}).

\section{The program and the message}

\subsection{How the program prints its results}

The program is one notebook cell. Its first line defines a helper, \code{show}, which prints
\textit{label = value} and writes integers in full, so even 617-digit numbers are printed digit by
digit. Each section below shows the lines of the program for one task and the lines they print.

\lstinputlisting[style=wl,firstline=1,lastline=1,firstnumber=1,caption={The helper that prints every intermediate value}]{../lab5.wl}

\subsection{From text to a number}

RSA and ElGamal work on integers, so the message must first become a number. As the handout says,
each character is replaced by its ASCII code in hexadecimal, the hexadecimal digits are written
one after another, and the result is read as one hexadecimal number and converted to decimal
(Table~\ref{tab:ascii}).

\begin{table}[H]
\centering
\caption{The message in ASCII}\label{tab:ascii}
\setlength{\tabcolsep}{3.3pt}
\small
\begin{tabular}{l*{15}{c}}
\toprule
Character & N & i & c & o & r & i & c & i & \textvisiblespace & D & a & n & i & e & l \\
Decimal & 78 & 105 & 99 & 111 & 114 & 105 & 99 & 105 & 32 & 68 & 97 & 110 & 105 & 101 & 108 \\
Hex & 4E & 69 & 63 & 6F & 72 & 69 & 63 & 69 & 20 & 44 & 61 & 6E & 69 & 65 & 6C \\
\bottomrule
\end{tabular}
\end{table}

\lstinputlisting[style=wl,firstline=3,lastline=11,firstnumber=3,caption={Converting the message to a number}]{../lab5.wl}
\lstinputlisting[caption={The message as a number}]{out_msg.txt}

\code{ToCharacterCode} gives the ASCII codes, \code{IntegerString[\dots, 16, 2]} writes each code as
two hexadecimal digits, and \code{FromDigits[hex, 16]} reads the joined string as a base-16 number.
The result, $m = 407136686053340208478179658767623532$, has 119 bits, far fewer than the 2048 bits
of the moduli below, as both algorithms require ($m < n$ for RSA and $m < p$ for ElGamal).
Decryption reverses the conversion with \code{IntegerDigits[m, 256]}, which splits the number back
into bytes, and \code{FromCharacterCode}, which turns the bytes back into characters.

\section{RSA (task 2.1)}

\subsection{The algorithm}

\textbf{Key generation.}
\begin{enumerate}
  \item Choose two different random primes $p$ and $q$ of 1024 bits each. The product of two
  1024-bit numbers has 2047 or 2048 bits, so the program repeats the choice until $n = pq$ has
  exactly 2048 bits, as the task requires.
  \item Compute the modulus $n = p \cdot q$ and Euler's function $\varphi(n) = (p - 1)(q - 1)$.
  \item Choose the public exponent $e$ with $\gcd(e, \varphi(n)) = 1$. I use $e = 65537 = 2^{16} + 1$,
  the value used in practice: it is prime and has only two 1-bits, so encryption is fast. The
  program checks the gcd together with $p$ and $q$ and chooses new primes if it is not 1.
  \item Compute the private exponent $d = e^{-1} \bmod \varphi(n)$, the number with
  $e \cdot d \equiv 1 \pmod{\varphi(n)}$. \code{PowerMod[e, -1, phi]} finds it with the extended
  Euclidean algorithm.
\end{enumerate}
The public key is $(n, e)$. The private key is $d$; $p$, $q$ and $\varphi(n)$ must also stay
secret, because anyone who knows them can compute $d$.

\textbf{Encryption} (by anyone, with the public key): $c = m^{e} \bmod n$.

\textbf{Decryption} (by the owner, with the private key): $m = c^{d} \bmod n$.

\textbf{Why it works.} Since $e d = 1 + t\,\varphi(n)$ for some integer $t$, Euler's theorem gives
$c^{d} = m^{e d} = m \cdot (m^{\varphi(n)})^{t} \equiv m \pmod{n}$.

\subsection{Computation}

\lstinputlisting[style=wl,firstline=13,lastline=34,firstnumber=13,caption={RSA in Wolfram Language}]{../lab5.wl}
\lstinputlisting[caption={RSA keys, cryptogram and decryption}]{out_rsa.txt}

\textbf{Observations:}
\begin{enumerate}
  \item \textbf{The decryption gives back the name.} \code{rsa.decrypted} equals $m$ and
  \code{rsa.text} is \textit{Nicorici Daniel}. The checks $\gcd(e, \varphi) = 1$ and
  $e \cdot d \bmod \varphi = 1$ both print 1, and $n$ has 2048 bits.
  \item \textbf{The cryptogram is as long as $n$.} The 119-bit message becomes a number of up to
  2048 bits, because $c$ can be any value below $n$.
  \item \textbf{Security rests on factoring.} $n$ and $e$ are public. Computing $d$ needs
  $\varphi(n)$, which needs $p$ and $q$, and nobody knows how to factor a 2048-bit $n$ in reasonable
  time.
  \item \textbf{Textbook RSA is deterministic.} Encrypting the same $m$ with the same key always
  gives the same $c$, so an attacker can test guesses. Real systems add random padding (OAEP) to
  the message before encrypting.
\end{enumerate}

\section{ElGamal (task 2.2)}

\subsection{The public parameters}

ElGamal works in the group of integers modulo a prime. The handout gives the prime $p$ and the
generator $g = 2$, which every user shares. The 617 digits of $p$ are written as eight strings of
80 digits only so that the code fits on the screen; \code{StringJoin} glues them together and
\code{ToExpression} reads the result as one integer.

\lstinputlisting[style=wl,firstline=36,lastline=48,firstnumber=36,caption={The public parameters from the handout}]{../lab5.wl}
\lstinputlisting[caption={The parameters as printed}]{out_params.txt}

The given $p$ is not an arbitrary prime: its hexadecimal form begins with
\code{FFFFFFFFFFFFFFFFADF85458A2BB4A9A} and ends with \code{FFFFFFFFFFFFFFFF}, which is the 2048-bit
group \code{ffdhe2048} standardised in RFC 7919 for Diffie-Hellman. It is a \textit{safe prime}:
$q = (p - 1)/2$ is also prime, which I checked in section~\ref{sec:verify}. A safe prime makes the
discrete logarithm problem as hard as possible for a modulus of its size.

\subsection{The algorithm}

\textbf{Key generation} (Bob):
\begin{enumerate}
  \item Choose a random private key $x$ with $1 < x < p - 1$.
  \item Compute the public key $\beta = g^{x} \bmod p$ and publish $(p, g, \beta)$.
\end{enumerate}

\textbf{Encryption} (Alice, with Bob's public key):
\begin{enumerate}
  \item Choose a random one-time number $k$ with $1 < k < p - 1$.
  \item Compute $y_1 = g^{k} \bmod p$ and $y_2 = m \cdot \beta^{k} \bmod p$.
  \item Send the pair $(y_1, y_2)$.
\end{enumerate}

\textbf{Decryption} (Bob, with $x$): $m = y_2 \cdot (y_1^{x})^{-1} \bmod p = y_2 \cdot y_1^{\,p-1-x} \bmod p$.

\textbf{Why it works.} $y_1^{x} = g^{k x} = (g^{x})^{k} = \beta^{k}$, so dividing $y_2$ by
$y_1^{x}$ removes the mask $\beta^{k}$ and leaves $m$. The inverse of $y_1^{x}$ is $y_1^{\,p-1-x}$,
because $y_1^{\,p-1} \equiv 1 \pmod p$ by Fermat's little theorem. This avoids computing a modular
inverse separately.

\subsection{Computation}

\lstinputlisting[style=wl,firstline=50,lastline=62,firstnumber=50,caption={ElGamal in Wolfram Language}]{../lab5.wl}
\lstinputlisting[caption={ElGamal keys, cryptogram and decryption}]{out_elgamal.txt}

\textbf{Observations:}
\begin{enumerate}
  \item \textbf{The decryption gives back the name.} \code{elgamal.decrypted} equals $m$ and
  \code{elgamal.text} is \textit{Nicorici Daniel}.
  \item \textbf{The cryptogram is twice as long as the message block.} It is the pair
  $(y_1, y_2)$, two numbers below $p$, about 4096 bits in total.
  \item \textbf{ElGamal is randomised.} Because $k$ is new for every message, encrypting the same
  name twice gives two different cryptograms. This is the property textbook RSA lacks.
  \item \textbf{$k$ must never be reused.} If two messages are encrypted with the same $k$, they
  share the same mask $\beta^{k}$, and knowing one message reveals the other:
  $m_2 = m_1 \cdot y_2' \cdot y_2^{-1} \bmod p$.
  \item \textbf{Security rests on the discrete logarithm.} Finding $x$ from $\beta = g^{x} \bmod p$
  is the discrete logarithm problem, which is infeasible for a 2048-bit safe prime.
\end{enumerate}

\section{Diffie-Hellman key exchange for AES-256 (task 3)}

\subsection{The algorithm}

Alice and Bob want to use AES-256, a symmetric cipher, so they need the same secret 256-bit key,
but they can only talk over an open channel. Diffie-Hellman lets them build that key in public.
Both use the public $p$ and $g = 2$ from the handout.

\begin{enumerate}
  \item Alice chooses a random secret $a$ and Bob a random secret $b$. As the algorithm requires,
  both are drawn uniformly from $2, \dots, p - 2$: the values $0$, $1$ and $p - 1$ are excluded
  because they would give the trivial public values $1$, $g$ and $\pm 1$.
  \item Alice computes $A = g^{a} \bmod p$ and sends it to Bob. Bob computes $B = g^{b} \bmod p$
  and sends it to Alice.
  \item Alice computes $K = B^{a} \bmod p$ and Bob computes $K = A^{b} \bmod p$. Both get the same
  number, because $B^{a} = g^{b a} = g^{a b} = A^{b}$.
  \item AES-256 needs exactly 256 bits, while $K$ is a 2048-bit number. Both sides take the 256
  most significant bits of the 2048-bit representation of $K$:
  $K_{\mathrm{AES}} = \lfloor K / 2^{1792} \rfloor$, computed with \code{Quotient}.
\end{enumerate}

\subsection{Computation}

\lstinputlisting[style=wl,firstline=64,lastline=80,firstnumber=64,caption={Diffie-Hellman and the AES key in Wolfram Language}]{../lab5.wl}
\lstinputlisting[caption={The exchange and the shared AES-256 key}]{out_dh.txt}

\textbf{Observations:}
\begin{enumerate}
  \item \textbf{Both sides get the same key.} \code{dh.keyAlice} and \code{dh.keyBob} are equal and
  \code{dh.equal} prints \code{True}, although $a$ and $b$ never left their owners.
  \item \textbf{An eavesdropper sees $p$, $g$, $A$ and $B$.} To find $K$ she would need $a$ or $b$,
  which means solving a discrete logarithm modulo a 2048-bit prime.
  \item \textbf{The AES key has exactly 256 bits.} It is printed as 64 hexadecimal digits and as
  256 binary digits. Leading zeros are kept, so the key length does not depend on the value.
  \item \textbf{In practice a key derivation function is used.} Real protocols such as TLS do not
  cut $K$ directly; they pass it through a key derivation function (for example HKDF with SHA-256),
  which spreads the randomness of all 2048 bits over the 256-bit key.
  \item \textbf{Plain Diffie-Hellman does not authenticate.} An attacker in the middle can run one
  exchange with Alice and another with Bob and read everything. Real protocols sign $A$ and $B$
  so each side knows who it is talking to.
\end{enumerate}

\section{Verification}\label{sec:verify}

To make sure every printed value is right, a Python script reads the output of the program and
recomputes each step on its own, with Python's built-in big integers: it tests that $p$ and $q$
are prime (Miller-Rabin), that $n = pq$ has 2048 bits, that $e d \equiv 1 \pmod{\varphi(n)}$, that
both encryptions and decryptions give the printed values and recover the name, that the
Diffie-Hellman keys agree, and that the AES key is the top 256 bits of the shared key. It also
confirms that the handout's $p$ is a safe prime.

\lstinputlisting[style=py,caption={verify.py},label={fig:verify}]{../verify.py}
\lstinputlisting[caption={All checks pass}]{verify_run.txt}

\plainsection{Conclusions}

I generated a 2048-bit RSA key pair, encrypted my name with the public key and decrypted it with the
private key. I did the same with ElGamal over the 2048-bit prime from the handout, and I ran a
Diffie-Hellman exchange between Alice and Bob that ended with the same 256-bit AES key on both
sides. Every value was checked independently.

The three algorithms rest on the same tool, modular exponentiation, which is easy to compute and
hard to reverse. RSA is hard to break because factoring $n$ is hard. ElGamal and Diffie-Hellman are
hard to break because computing discrete logarithms is hard.

The lab also showed why these algorithms are used the way they are. Public-key operations on
2048-bit numbers are slow and their cryptograms are large, so in practice they are not used to
encrypt data directly. They are used to agree on or transport a short symmetric key, exactly as in
task 3, and the data itself is then encrypted with a fast symmetric cipher such as AES.
\end{document}
